
\chapter{Introducción}\label{cap:introduccion}

En la actualidad, muchas empresas dependen de servicios digitales para asegurar la continuidad operativa y la protección de sus transacciones \cite{rashid2025establishing}. En este entorno, la Infraestructura de Clave Pública (PKI) se ha consolidado como el pilar tecnológico que sostiene la seguridad y la confianza en el ecosistema digital global, constituyéndose en la base de las transacciones electrónicas y del negocio electrónico \cite{prodanovic2019_pki, rashid2025establishing}. Su propósito principal es asegurar la confiabilidad de las interacciones digitales, estableciendo mecanismos que garanticen autenticidad, integridad y confidencialidad \cite{rashid2025establishing}. Dentro de este marco, las Autoridades de Certificación (CA) cumplen un rol esencial, actuando como entidades responsables de verificar identidades y gestionar el ciclo de vida de los certificados digitales, los cuales vinculan una identidad con una clave criptográfica \cite{mana2019_swimpki}. Estos certificados se convierten en elementos indispensables para habilitar transacciones seguras en diferentes sectores \cite{rashid2025establishing}.

En el ámbito de las CA, los sistemas tradicionales han sido analizados debido a distintas deficiencias que afectan su funcionamiento y evolución \cite{wu2009_sshca}. Una de las más relevantes, especialmente en arquitecturas jerárquicas o simples, es la rigidez generada por implementaciones fuertemente acopladas \cite{wu2009_sshca}. Este acoplamiento limita la expansión de servicios y dificulta la incorporación de nuevos servicios, restringiendo tanto la eficiencia como la calidad del desarrollo de software \cite{wu2009_sshca}. Aunque el modelo jerárquico permite una expansión inicial adecuada, concentra un punto crítico de vulnerabilidad al depender completamente de la Root CA, lo que implica un riesgo alto si esta llega a comprometerse \cite{prodanovic2019_pki}. A esto se añade que muchos sistemas heredados mantienen estructuras poco flexibles, lo cual dificulta su modernización y la introducción de mejoras \cite{wu2009_sshca}.

Otra deficiencia identificada se relaciona con la centralización de los procesos de validación, lo que deja a los usuarios expuestos a ataques tales como Man-In-The-Middle (MITM) en situaciones donde los servicios de verificación pueden ser manipulados \cite{yasin2019_korgan}. Asimismo, los sistemas CA y PKI presentan dificultades para manejar un creciente volumen de usuarios y transacciones, lo que pone a prueba la capacidad de respuesta de las plataformas tradicionales \cite{wu2009_sshca}. Estas limitaciones estructurales dificultan que una CA evolucione al ritmo que demandan las nuevas necesidades tecnológicas y los requisitos de seguridad contemporáneos \cite{Rezaei2025}.

Ante las deficiencias identificadas como el acoplamiento fuerte entre componentes, la dificultad para incorporar nuevos servicios, las vulnerabilidades asociadas a puntos únicos de falla que facilitan ataques como MITM y la capacidad limitada para manejar altos volúmenes de usuarios y transacciones, se propone el diseño de una Arquitectura de Software capaz de superar las restricciones de los modelos tradicionales y ofrecer una estructura más flexible, modular y resiliente. La propuesta busca aplicar principios de separación estricta de funciones esenciales a través de componentes pequeños e independientes que permitan ampliar servicios sin comprometer la estabilidad general del sistema. Además, con el fin de reducir la dependencia de componentes centrales y mitigar los riesgos asociados a puntos únicos de falla, pueden considerarse enfoques descentralizados, entre ellos modelos basados en Practical Byzantine Fault Tolerance (PBFT), que permiten eliminar la dependencia de entidades intermedias para la validación y distribuir la responsabilidad operativa dentro de la arquitectura \cite{yasin2019_korgan}. En paralelo, el manejo de claves privadas se apoyará en módulos criptográficos de hardware (HSM) que garanticen niveles adecuados de protección y cumplimiento normativo \cite{aciobanitei2020_learned, chesnokov2020_eds}. Finalmente, la integración de tecnología blockchain se incorpora como un elemento arquitectónico relevante para asegurar inmutabilidad y trazabilidad precisa en los procesos de auditoría \cite{Rezaei2025}.

El objetivo principal de este diseño es establecer una arquitectura de referencia que pueda ser adoptada por diferentes Autoridades de Certificación que operan bajo el marco legal ecuatoriano y que requieren estructuras más flexibles que faciliten la expansión de servicios y la reutilización de componentes. La investigación examinará diferentes enfoques arquitectónicos como los modelos modulares, orientados a servicios y basados en eventos, con el fin de combinar sus principales ventajas y construir una arquitectura equilibrada en términos de seguridad, mantenibilidad y capacidad de evolución \cite{prodanovic2019_pki}.

\subsection*{Objetivos de Investigación}

El objetivo general del presente trabajo es desarrollar un prototipo de arquitectura
de referencia para Autoridades de Certificación reguladas bajo el marco legal ecuatoriano,
enfocada en la modularidad y escalabilidad, que sirva como modelo para superar las
limitaciones de rigidez y acoplamiento de los sistemas PKI tradicionales.

Para alcanzar este objetivo general se definen los siguientes cuatro objetivos específicos.
\begin{enumerate}
  \item Analizar los estándares PKI (X.509, OCSP), los modelos de confianza y los marcos
        regulatorios nacionales (ARCOTEL y MINTEL) e internacionales relevantes (NIST/eIDAS),
        para definir los requisitos funcionales y no funcionales que debe cumplir una
        arquitectura de CA moderna en Ecuador.
  \item Diseñar la arquitectura de referencia focalizada en los componentes críticos de la CA
        (emisión, gestión, validación y auditoría), aplicando patrones de modularidad y
        seguridad, y formalizando la estructura mediante la elaboración de artefactos técnicos
        tangibles usando modelos C4 bajo la metodología TOGAF.
  \item Implementar una prueba de concepto (PoC) del módulo de auditoría en un entorno de
        simulación, integrando un componente de tecnología blockchain (Hyperledger Fabric)
        para verificar la viabilidad y eficacia de los principios definidos en la arquitectura
        de referencia.
  \item Evaluar el diseño de la arquitectura de referencia mediante el método Architecture
        Tradeoff Analysis Method (ATAM) para identificar riesgos y compromisos en atributos
        de calidad; y validar el desempeño del prototipo desarrollado mediante métricas clave
        (latencia de emisión, acoplamiento eferente Ce).
\end{enumerate}

\subsection*{Organización del Documento}
El presente documento se organiza en los siguientes capítulos. El Capítulo 2 (Antecedentes) contextualiza el estado actual de la PKI, describiendo la evolución de sus arquitecturas, los problemas estructurales de los sistemas tradicionales y los enfoques emergentes que motivan esta investigación. El Capítulo 3 (Marco Teórico) establece los fundamentos conceptuales, normativos y metodológicos del trabajo: los principios criptográficos y el estándar X.509, el marco legal ecuatoriano encabezado por la Resolución ARCOTEL-2024-0176, los estilos y patrones de arquitectura de software, y los \textit{frameworks} de diseño y evaluación TOGAF ADM y ATAM. El Capítulo 4 (Trabajos Relacionados) revisa las propuestas existentes en la literatura sobre arquitecturas PKI modulares y descentralizadas, identificando brechas existentes. El Capítulo 5 (Metodología) constituye el núcleo del trabajo: detalla las cuatro fases del proceso de diseño, análisis de requisitos normativos y funcionales, construcción del modelo arquitectónico mediante TOGAF ADM (fases Preliminary a D), diseño e implementación del módulo de auditoría, y validación de la arquitectura mediante ATAM, junto con los artefactos producidos en cada etapa, las decisiones arquitectónicas documentadas en ADRs y los escenarios del Árbol de Utilidad evaluados. El Capítulo 6 (Resultados) presenta la arquitectura de referencia obtenida, los diagramas C4, las métricas de acoplamiento y los resultados de las pruebas del prototipo de auditoría. Finalmente, el Capítulo 7 (Conclusiones) sintetiza los hallazgos, discute las implicaciones de las decisiones arquitectónicas para las CA ecuatorianas, enumera las limitaciones del estudio y propone líneas de trabajo futuro.
